\documentclass[12pt]{article}
\usepackage{amsmath,amssymb,amsthm}
\usepackage[margin=1in]{geometry}
\usepackage{array}
\usepackage{hyperref}
\usepackage{xurl}
\emergencystretch=3em
% Theorem environments are declared inside
% sec-carry.tex (first \input below); the
% section fragments guard their own declarations with
% \@ifundefined and also compile standalone.
% NOTE (deposit copy): the source tree keeps the section
% fragments under sections/; in this flat archive they sit
% next to this master file and the \input lines below were
% adjusted accordingly (sections/sec-X -> sec-X).

\title{Palindromes $p$ with $kp+1$ palindromic:
carry-chain classifications for $2 \le k \le 6$ and
certified zeros of OEIS A083837}
\author{Fr\'ed\'eric Gaudin\thanks{Universit\'e de Caen
Normandie (UNICAEN), Caen, France. Email:
\texttt{frederic.gaudin@unicaen.fr}}%
\thanks{2020 Mathematics Subject Classification: 11A63,
11B83, 11Y99.}%
\thanks{Keywords: palindrome, decimal digits, carry
propagation, repdigit, repunit, OEIS A069881, A083829,
A083830, A083831, A083832, A083833, A083834, A083835,
A083836, A083837, A083838, formal verification, Lean.}}
\date{November 2026}

\begin{document}
\maketitle

\begin{abstract}
For $k \in \{2,3,4,5,6\}$ let $S_k$ be the set of decimal
palindromes $p$ such that $kp+1$ is again a palindrome:
OEIS A069881 ($k=2$, Tangente problem 8154), A083829
($k=3$), A083831 ($k=4$), A083833 ($k=5$) and A083835
($k=6$), with the image sequences A083830, A083832,
A083834 and A083836. Every published conjecture on these
entries -- the \emph{Tangente} problem and the recorded
observations of R.~Israel and M.~S.~Branicky -- is proved
here and certified in the Lean~4
kernel: the even-length members are exactly the repdigits
($55\ldots5$, $77\ldots7$, $88\ldots8$, $22\ldots2$,
$99\ldots9$ respectively), each an iff in both
directions, the digit alphabets are the published bounds,
and the $k=6$ entry is settled by a complete
classification. The same carry-chain device -- a chain
folded back onto itself by palindromicity, $2$-periodic
at constant length and frozen by any overflow -- yields,
at the other end of the family, the twenty-one displayed
zeros of A083837: each is upgraded from Chandler's 2003
finite-search bound ($< 10^{15}$) to a universal
nonexistence theorem by a prefix-suffix obstruction, with
the twenty-one zeros of A083838. All results rest on the
three standard axioms; no {\ttfamily native\_decide} is
used.
\end{abstract}

\section{Introduction}
\label{sec:intro}

For a positive integer $k$, write $S_k$ for the set of
decimal palindromes $p$ such that $q = kp+1$ is again a
palindrome. The On-Line Encyclopedia of Integer
Sequences~\cite{oeis} holds one entry per multiplier
$k = 2, \dots, 6$ together with, for $k \ge 3$, the image
entry listing the values $q$; all six source entries were
contributed by A.~Murthy and M.~Srikanth in May 2003
(A069881: Murthy, April 2002) and completed by
R.~Chandler the same month. The definitions read
verbatim:
\begin{quote}
A069881: ``Numbers $n$ such that $n$ and $2n+1$ are both
palindromes.''\\
A083829: ``Palindromes $k$ such that $3k + 1$ is also a
palindrome.''\\
A083831: ``Palindromes $n$ such that $4n + 1$ is also a
palindrome.''\\
A083833: ``Palindromes $p$ such that $5p + 1$ is also a
palindrome.''\\
A083835: ``Palindromes $n$ such that $6n + 1$ is also a
palindrome.''
\end{quote}
The image entries A083830, A083832, A083834 and A083836
list the ``palindromes arising in'' A083829, A083831,
A083833 and A083835 respectively. Two further entries
scan the family over the multiplier:
\begin{quote}
A083837: ``Smallest palindrome $n$ such that $kn + 1$
is a palindrome, or $0$ if no such number exists.''\\
A083838: ``Smallest palindrome of the form $nk+1$ where
$k$ is also a palindrome:, or $0$ if no such number
exists; palindromes arising in A083837.''
\end{quote}
(the trailing colon in the A083838 definition is the
entry's own). On A083837 and A083838 the same comment
appears, dated July 13, 2003: ``0 in sequence indicates
no palindrome $< 10^{15}$'' (R.~Chandler): a displayed
zero is a finite-search observation, not a theorem of
nonexistence.

\paragraph{Published conjectures on the entries.}
Four dated conjectural items are attached to these
entries. On A069881, the classification of even-length
members is the subject of \emph{Tangente}
problem~8154~\cite{tangente}, posed after J.~McCann's
2002 comment that every repdigit $55\ldots5$ is a term
(``on peut conjecturer que les seuls palindromes
$\alpha+$ ayant un nombre pair de chiffres sont $55$,
$5555$, $555555$, \ldots''). On A083831, R.~Israel
recorded (April 4, 2018): ``Among infinite subsequences
are the repdigits $8\ldots8 = 8(10^k-1)/9$. It appears
that the only terms with an even number of digits are
these for even $k$.'' On A083829 he recorded (February
23, 2023, tagged ``It appears'') four observations:
the only even-length terms are the repdigits of
A002281; all other terms of more than one digit start
with $14$, $15$, $16$, $24$, $25$ or $26$; no term
contains the digits $3$, $8$ or $9$; and the only terms
containing $7$ are those of A002281. On A083833 and
A083834, M.~S.~Branicky recorded (June 13, 2021,
tagged ``(Conjectures)'') the digit alphabets
$\{0,1,2,3\}$ and $\{0,1,5,6\}$ and the parity claims
that even-length terms are the repdigit $22\ldots2$ and
the repunit $11\ldots1$; the A083834 verbatim contains
an even/odd slip corrected in
Remark~\ref{rem:parity-slip}.

\paragraph{Contents.}
Section~\ref{sec:carry} presents the shared carry-chain
device once, for every $k \in \{2,\dots,6\}$: at constant
length the chain is symmetric and $2$-periodic, forcing
odd length and alternating digit classes; at overflow the
chain freezes on a single repdigit family selected by the
head equation $kd+1 = 11v$. The instantiations settle
all the conjectures above -- $k=2$, A069881
(Section~\ref{sec:k2}); $k=3$, A083829 and A083830
(Section~\ref{sec:k3}); $k=4$, A083831 and A083832
(Section~\ref{sec:k4}); $k=5$, A083833 and A083834
(Section~\ref{sec:k5}); $k=6$, A083835 and A083836
(Section~\ref{sec:k6}) -- each even-length classification
as an iff in both directions. The twenty-one displayed
zeros of A083837 (Theorem~\ref{thm:zeros}) and of A083838
(Corollary~\ref{cor:zeros838}) are proved in
Section~\ref{sec:zeros} by a prefix-suffix obstruction,
turning Chandler's bound into universal nonexistence.
Section~\ref{sec:lean} records exactly what is
kernel-verified in Lean~4 and in which modules, with the
paper--kernel correspondence of Table~\ref{tab:lean}.

\paragraph{Related entries, cited only.}
Two additive variants sit next to this family and are
already settled on their own entries: on A082274
(palindromes $p$ with $p+2$ palindromic), F.~Fr\"ohlich's
2014 argument on the entry shows the members are
$\{1,\dots,7,9\}$ together with the repdigits
$99\ldots9$; on A082275 ($p+11$ palindromic), D.~A.~Corneth's
2025 comment proves the members are exactly the
palindromes $d\,9^{t}\,d$ with $d \in \{1,\dots,8\}$
\cite{a082274,a082275}. The additive shift kills the
interior carry freedom entirely, whereas multipliers keep
the alternating classes of Section~\ref{sec:carry}; the
two mechanisms do not overlap and these entries are only
pointed to, not absorbed. The repdigit and repunit
families occurring throughout are OEIS A002279, A002113
and A000042. In the wider literature on decimal
palindromes we point to Harminc and Sot\'ak's study of
palindromic numbers in arithmetic
progressions~\cite{harmincsotak}, the theorem of Banks,
Hart and Sakata that almost all palindromes are
composite~\cite{banks}, and the three-palindromes
additive basis of Cilleruelo, Luca and
Baxter~\cite{cll}.

% sec-carry.tex -- the shared carry-chain device for the
% kp+1 palindrome developments (k in {2,3,4,5,6}).
% Standalone fragment, \input by the master file.
% The theorem environments are declared here (first input
% in document order); the later section files guard their
% own declarations with \@ifundefined so that every
% fragment also compiles standalone.

\theoremstyle{plain}
\newtheorem{theorem}{Theorem}
\newtheorem{lemma}[theorem]{Lemma}
\newtheorem{corollary}[theorem]{Corollary}
\newtheorem{question}[theorem]{Question}
\theoremstyle{remark}
\newtheorem{remark}[theorem]{Remark}
\providecommand{\IsPal}{\mathrm{IsPal}}

\section{The shared carry-chain device}
\label{sec:carry}

Fix $k \in \{2,3,4,5,6\}$. Let $p$ be a decimal
palindrome of $m \ge 1$ digits, and assume $q = kp+1$ is
again a palindrome, of $m'$ digits. Write
$a_j = \lfloor p/10^j \rfloor \bmod 10$ for the $j$-th
digit of $p$ (units $= j = 0$) and define the carry into
position $j$ of the computation of $kp+1$ (the usual
base-$10$ addition carry, cf.~\cite{holte}) by the closed
form
\begin{gather*}
 c_j \;=\; \bigl\lfloor\bigl(k\,(p \bmod 10^j) + 1\bigr)\,
 \big/\,10^j\bigr\rfloor ,
 \\
 \text{so that}\qquad
 b_j = (k a_j + c_j) \bmod 10 ,
 \qquad c_{j+1} = \bigl\lfloor\bigl(k a_j + c_j\bigr)\big/10
 \bigr\rfloor ,
\end{gather*}
where $b_j$ is the $j$-th digit of $q$. Then $c_0 = 1$
(the $+1$ is folded into the chain), $0 \le c_j \le k-1$,
and since $kp+1 < 10^{m+1}$ the length satisfies
$m' \in \{m, m+1\}$, with $m' = m$ iff $c_m = 0$.
Palindromicity of $q$ reads $b_j = b_{m'-1-j}$. We write
$R_m = (10^m-1)/9$ for the $m$-digit repunit and
$d\,R_m$ for the repdigit $dd\dots d$. We write
$\IsPal(p,m)$ for the predicate that $p$ is a positive
$m$-digit integer ($1 \le p \le 9$ for $m = 1$, else
$10^{m-1} \le p < 10^m$) whose base-$10$ digits satisfy
$a_j = a_{m-1-j}$ for all $0 \le j < m$.

\begin{lemma}[same length: symmetric $2$-periodic chain]
\label{lem:same}
If $m' = m$ then, for every $k \in \{2,3,4,5,6\}$:
\emph{(i)} $c_j = c_{m-1-j}$ for all $j < m$;
\emph{(ii)} $c_{j+1} = c_{j-1}$ for $1 \le j \le m-2$,
hence $c_{2t} = 1$ for $2t \le m-1$ and
$c_{2t+1} = c_1$ for $2t+1 \le m-1$;
\emph{(iii)} $m$ is odd and $c_1 = 0$;
\emph{(iv)} the digits obey the alternating classes
$a_j \le \lfloor 8/k \rfloor$ at even positions and
$\lceil 10/k \rceil \le a_j \le \lfloor 19/k \rfloor$ at
odd positions; in particular
$a_0 \in \{1,\dots,\lfloor 8/k \rfloor\}$.
\end{lemma}

\begin{proof}
(i) $b_j = b_{m-1-j}$ reads
$(k a_j + c_j) \equiv (k a_{m-1-j} + c_{m-1-j})
\pmod{10}$; the input symmetry $a_{m-1-j} = a_j$ cancels
$k a_j$ exactly, and the residues
$c_j, c_{m-1-j} \le k-1 < 10$ coincide. (ii) Apply the
recurrence at $j$ and at the mirror index $m-1-j$, then
rewrite through (i) and the same digit symmetry. (iii) If
$m$ is even, $m-1$ is odd: (i) at $j = 0$ gives
$c_{m-1} = c_0 = 1$ while (ii) gives $c_{m-1} = c_1$, so
the top recurrence reads
$c_m = \lfloor(k a_0 + c_{m-1})/10\rfloor
= \lfloor(k a_0 + 1)/10\rfloor = c_1 = 1$,
contradicting $c_m = 0$. Hence $m$ is odd, $m-1$ is
even, $c_{m-1} = 1$ by (ii), and
$0 = c_m = \lfloor(k a_0 + 1)/10\rfloor$ forces
$k a_0 \le 8$ and $c_1 = 0$. (iv) At even $j$ one has
$c_j = 1$ and $c_{j+1} = 0$ (odd index, or $j+1 = m$), so
$\lfloor(k a_j + 1)/10\rfloor = 0$, i.e.\ $k a_j \le 8$.
At odd $j$, $c_j = c_1 = 0$ and $c_{j+1} = 1$ (even
index), so $10 \le k a_j \le 19$.
\end{proof}

Table~\ref{tab:classes} lists the alternating classes of
Lemma~\ref{lem:same}(iv) for each multiplier.

\begin{table}[ht]
\centering
\small
\begin{tabular}{c|c|c|c|c}
\hline
$k$ & entry & even positions & odd positions &
$a_0$ (same length) \\
\hline
2 & A069881 & $\{0,1,2,3,4\}$ & $\{5,6,7,8,9\}$ &
$\{1,2,3,4\}$ \\
3 & A083829 & $\{0,1,2\}$ & $\{4,5,6\}$ & $\{1,2\}$ \\
4 & A083831 & $\{0,1,2\}$ & $\{3,4\}$ & $\{1,2\}$ \\
5 & A083833 & $\{0,1\}$ & $\{2,3\}$ & $\{1\}$ \\
6 & A083835 & $\{0,1\}$ & $\{2,3\}$ & $\{1\}$ \\
\hline
\end{tabular}
\caption{Digit classes forced on same-length members, by
multiplier. Even-length members are excluded from this
table: they all live in the overflow branch
(Lemma~\ref{lem:overflow}).}
\label{tab:classes}
\end{table}

\begin{lemma}[sufficiency of the alternating classes]
\label{lem:altsuff}
Let $k \in \{2,3,4,5,6\}$, let $m$ be odd, and let $p$ be
an $m$-digit palindrome whose digits belong to the
alternating classes of Table~\ref{tab:classes}: $a_0 \ge
1$, $k a_j \le 8$ at even positions and
$10 \le k a_j \le 19$ at odd positions. Then $kp+1$ is an
$m$-digit palindrome.
\end{lemma}

\begin{proof}
The chain closes by induction: $c_0 = 1$; at an even
position $k a_{2t} \le 8$ gives $b_{2t} = k a_{2t} + 1 \le
9$ and $c_{2t+1} = 0$; at an odd position
$10 \le k a_{2t+1} \le 19$ gives $b_{2t+1} = k a_{2t+1} -
10$ and $c_{2t+2} = 1$. Hence $c_{2t} = 1$ and
$c_{2t+1} = 0$ throughout; as $m-1$ is even the top carry
is $c_m = 0$, so $q = kp+1$ keeps length $m$. Since $j$
and $m-1-j$ have the same parity and the output digit
depends only on the parity of the position
($b_{2t} = k a_{2t} + 1$, $b_{2t+1} = k a_{2t+1} - 10$),
$a_{m-1-j} = a_j$ gives $b_{m-1-j} = b_j$.
\end{proof}

\begin{lemma}[overflow: the frozen chain]
\label{lem:overflow}
If $m' = m+1$ the head equation
\begin{equation}
 (k a_0 + 1) \bmod 10 \;=\; c_m ,
 \qquad c_m = \bigl\lfloor\bigl(k a_0 + c_{m-1}\bigr)\big/10
 \bigr\rfloor
 \label{eq:head}
\end{equation}
(the first digit $b_0$ equals the new leading digit
$b_m = c_m$) has, for each $k \in \{2,3,4,5,6\}$, a single
solution: $a_0 = d_k$ and $c_m = v_k$, where $(d_k, v_k)$
is the unique digit solution of $k d + 1 = 11 v$, namely
\[
 (k;\,d_k;\,v_k) \;=\;
 (2;5;1),\; (3;7;2),\; (4;8;3),\; (5;2;1),\; (6;9;5) .
\]
The whole chain freezes: $c_j = v_k$ for $1 \le j \le m$
and every digit of $p$ equals $d_k$. Hence $p = d_k\,R_m$
is a repdigit, and $q = kp+1$ is the repunit $R_{m+1}$
for $k \in \{2,5\}$, respectively the framed repdigit
$v_k\,w_k^{\,m-1}\,v_k$ with
$w_k = (k d_k + v_k) \bmod 10$ for $k \in \{3,4,6\}$,
i.e.\ $2\,3^{\,m-1}\,2$, $3\,5^{\,m-1}\,3$ and
$5\,9^{\,m-1}\,5$.
\end{lemma}

\begin{proof}
\emph{Scheme I} ($k \in \{2,3,4,6\}$). The head equation
\eqref{eq:head} is solved by a finite check on
$a_0 \in \{1,\dots,9\}$ against $1 \le c_m \le k-1$ and
$c_m = \lfloor(k a_0 + c_{m-1})/10\rfloor$: only
$a_0 = d_k$ survives, giving $c_m = v_k$ and $c_1 = v_k$.
The propagation at position $j+1$ uses the mirror couple
\emph{(A)}
$(k a_{j+1} + c_{j+1}) \equiv (k a_j + c_{m-j-1})
\pmod{10}$ (the output symmetry folded by
$a_{m-j-1} = a_j$) and \emph{(B)}
$c_{m-j-1} = \lfloor(k a_{j+1} + c_{m-j-2})/10\rfloor$
(the mirrored recurrence). For $k \in \{2,4,6\}$ the pair
\emph{(A)}+\emph{(B)} has the single fixed point
$(a_{j+1}, c_{j+2}) = (d_k, v_k)$: the chain is rigid in
one step (for $k = 2$ the parity of $2a'+1$ forces the
mirror carry to $1$, then $a' \in \{0,5\}$ meets
$a' \ge 5$; for $k = 4$ and $k = 6$ the joint system is
solved digit by digit over the ten possibilities of
$a'$). For $k = 3$ the state $(a',c') = (3,1)$ is a
spurious fixed point of \emph{(A)}+\emph{(B)}; it is
killed one position later by the same equations shifted:
the shifted mirror congruence forces the next digit to
be $6$ (so $3 \cdot 6 = 18$), and the shifted mirrored
recurrence would then require the carry
$0 = \lfloor(18 + c)/10\rfloor$, impossible, so the induction runs
over $j \le m-2$ and the last digit is closed by
palindromy $a_{m-1} = a_0$.
\emph{Scheme II} ($k = 5$). The multiplier $5$ separates
carries mod $10$: $5a \bmod 10 \in \{0,5\}$ leaves
$c < 5$ readable in $b_j$, so the mirror congruence
yields directly $a_j \equiv a_{j-1} \pmod 2$ \emph{and}
$c_j = c_{m-j}$ ($1 \le j \le m-1$). A $2$-cycle
$g(u,a) = w$, $g(w,a) = u$ with $u \ne w$ is impossible
($11\,|u-w| \le 9$), hence the chain is constant
$c_j = v$ without induction; \eqref{eq:head} forces
$v = 1$ and $a_0 = 2$, and
$\lfloor(5 a_j + 1)/10\rfloor = 1$ with the common parity
gives $a_j = 2$ throughout. In all cases $b_0 = b_m = v_k$
and $b_j = (k d_k + v_k) \bmod 10$ inside, which is the
displayed image form.
\end{proof}

\begin{remark}[what the factorization actually shares]
\label{rem:honest}
The same-length half (Lemma~\ref{lem:same}) is literally
the same proof in all five certified developments; the
overflow half factors the \emph{statement} (frozen
constant chain, repdigit member) but not the proof:
$k \in \{2,4,6\}$ is rigid in one step, $k = 3$ needs the
one-step dead-end lookahead, and $k = 5$ goes through the
preserved mirror symmetry and the no-$2$-cycle argument,
a route closed to the other multipliers where
$k a \bmod 10$ does not separate the carry. In the kernel
development the machinery is likewise duplicated per
multiplier: there is no parametrized general lemma, the
same carry-chain argument is formalized for each $k$
(Section~\ref{sec:lean}).
\end{remark}

\begin{remark}[the multipliers $k = 1$ and
$k \in \{7,8,9\}$]
\label{rem:k1}
For $k = 1$ the class is finite: $S_1 = \{1,\dots,8\}$.
An $m$-digit palindrome with $m \ge 2$ and $a_0 \le 8$
would give $p+1$ the same leading digit $a_0$ but last
digit $a_0+1$, never palindromic; $a_0 = 9$ makes $p+1$
end in $0$, never a palindrome either. This elementary
case needs no machine. For $k \in \{7,8,9\}$ the head
key $kd+1 = 11v$ still has solutions -- $(d,v) = (3,2)$,
$(4,3)$, $(6,5)$ -- predicting repdigit overflow families
$33\ldots3 \mapsto 2\,3^{\,m-1}\,2$,
$44\ldots4 \mapsto 3\,5^{\,m-1}\,3$ and
$66\ldots6 \mapsto 5\,9^{\,m-1}\,5$; the displayed terms
of A083837 agree ($a(7) = a(8) = 1$, $a(9) = 6$, and
$7\cdot33+1 = 232$, $8\cdot44+1 = 353$,
$9\cdot66+1 = 595$ are palindromes). The full
classification of those entries is outside the certified
perimeter of this paper.
\end{remark}

% sec-A069881.tex -- k = 2 section of the unified kp+1 paper.
% Standalone fragment, \input by A083833-UNIFIE-paper.tex.
% Environments and notation inherited from sec-carry.tex; the
% guards below are no-ops once defined there.
\providecommand{\IsPal}{\mathrm{IsPal}}
\makeatletter
\@ifundefined{theorem}{\newtheorem{theorem}{Theorem}}{}
\@ifundefined{corollary}{\newtheorem{corollary}[theorem]{Corollary}}{}
\makeatother

\section{Multiplier $k = 2$: A069881}
\label{sec:k2}

Entry A069881~\cite{oeis} (A.~Murthy, April 2002) lists
$1, 2, 3, 4, 5, 55, 151, 161, 171, 181, 191, 252, \ldots$
Its only comment (J.~McCann, 2002) records that every
repdigit $55\ldots5$ is a term, and the classification
of even-length members is the conjecture of
\emph{Tangente} problem~8154~\cite{tangente} quoted in
Section~\ref{sec:intro}. Instantiating the carry chain
of Section~\ref{sec:carry} at $k = 2$ gives
$c_j = \bigl\lfloor\bigl(2\,(p \bmod 10^j) + 1\bigr)
/10^j\bigr\rfloor$, hence $c_0 = 1$ and $c_j \le 1$.

\begin{theorem}[Tangente 8154, full form]
\label{thm:k2-even}
Let $m \ge 1$ be even and $\IsPal(p,m)$. Then
$2p+1$ is a palindrome (of some length $m'$) if and
only if every digit of $p$ equals $5$, i.e.\ $p$ is the
repdigit $5\,R_m$.
\end{theorem}

\begin{theorem}[member dichotomy]
\label{thm:k2-dicho}
If $\IsPal(p,m)$ and $\IsPal(2p+1,m')$ with $m \ge 1$,
then either every digit of $p$ equals $5$ (the overflow
branch, where $2p+1 = 11\ldots1$ has $m+1$ digits), or
$m$ is odd, $a_0 \in \{1,2,3,4\}$, and $a_j \le 4$ at
even positions while $a_j \ge 5$ at odd positions.
\end{theorem}

\begin{corollary}[image parity]
\label{cor:k2-img}
Under the same hypotheses, every digit of $q = 2p+1$ at
an even position is odd.
\end{corollary}

\begin{proof}
Same-length branch ($m' = m$):
Lemma~\ref{lem:same} at $k = 2$ forces $m$ odd,
$a_0 \in \{1,2,3,4\}$ and the displayed classes.
Overflow branch ($m' = m+1$): the boundary
$(2a_0 + 1) \bmod 10 = c_m = 1$ forces $a_0 = 5$ and
$c_1 = 1$; at each subsequent position the mirrored
carry is forced to $1$ by parity ($2a'+1$ is odd),
while the joint constraints give $a' \in \{0,5\}$ and
$a' \ge 5$: the chain is rigid in a single step, so
$a_j = 5$ throughout, $p = 55\ldots5$ and
$q = 11\ldots1$. The converse direction of
Theorem~\ref{thm:k2-even} is the kernel-verified
identity $2 \cdot 5R_m + 1 = R_{m+1}$ packaged as a
literal iff (Table~\ref{tab:lean}); the forward
direction follows because even $m$ excludes the
same-length branch.
\end{proof}

The statements are consistent with all $45$ published
terms: the only even-length ones are the repdigits $55$
and $5555$, and every other listed term has odd length,
units digit in $\{1,2,3,4\}$ and the alternating classes
of Theorem~\ref{thm:k2-dicho}.

% sec-A083829.tex -- k = 3 section of the unified kp+1 paper.
% Standalone fragment, \input by A083833-UNIFIE-paper.tex.
% Environments and notation inherited from sec-carry.tex; the
% guards below are no-ops once defined there.
\providecommand{\IsPal}{\mathrm{IsPal}}
\makeatletter
\@ifundefined{theorem}{\newtheorem{theorem}{Theorem}}{}
\@ifundefined{corollary}{\newtheorem{corollary}[theorem]{Corollary}}{}
\@ifundefined{remark}{\newtheorem{remark}[theorem]{Remark}}{}
\makeatother

\section{Multiplier $k = 3$: A083829 and its image A083830}
\label{sec:k3}

On the entry A083829~\cite{oeis} -- palindromes $p$ such
that $q = 3p+1$ is again a palindrome -- Robert Israel
(February 23, 2023, tagged ``It appears'') recorded four
observations, numbered here P1--P4 and quoted in
Section~\ref{sec:intro}. The carry chain of
Section~\ref{sec:carry} at $k = 3$ has
$c_j = \bigl\lfloor\bigl(3\,(p \bmod 10^{j}) + 1\bigr)
/10^{j}\bigr\rfloor$, so $c_0 = 1$ and $c_j \le 2$.

\begin{theorem}[Israel's observations on A083829]
\label{thm:k3-israel}
Let $\IsPal(p,m)$ and $\IsPal(3p+1,m')$ with $m \ge 1$.
Then:
\begin{itemize}
\item[(P1)] if $m$ is even, every digit of $p$ is $7$
($p$ has the shape $77\ldots7$ of A002281);
\item[(P2)] if $m \ge 2$ and some digit of $p$ differs
from $7$ -- equivalently $p$ is not the all-$7$
repdigit -- the leading pair $10\,a_{m-1} + a_{m-2}$
lies in $\{14, 15, 16, 24, 25, 26\}$;
\item[(P3)] no digit of $p$ lies in $\{3, 8, 9\}$;
\item[(P4)] if some digit of $p$ equals $7$, then every
digit equals $7$.
\end{itemize}
\end{theorem}

\begin{corollary}[alphabet of the image; A083830]
\label{cor:k3-alphabet}
Under the same hypotheses, every digit of $q = 3p+1$
avoids $\{0, 6, 9\}$, i.e.\ lies in
$\{1,2,3,4,5,7,8\}$.
\end{corollary}

\begin{proof}
Two branches of the shared carry lemma. \emph{Same
length} ($m' = m$): Lemma~\ref{lem:same} forces a
$2$-periodic chain, $m$ odd, $a_0 \in \{1,2\}$ and
$c_1 = 0$, confining the digits to
$a_j \in \{0,1,2\}$ at even positions and
$a_j \in \{4,5,6\}$ at odd positions. \emph{Overflow}
($m' = m+1$): the head equation forces $a_0 = 7$ and
$c_1 = 2$; inductively the joint constraints admit one
spurious state $(a_{j+1}, c_{j+2}) = (3,1)$, locally
consistent but a dead end -- at the next position the
shifted mirror congruence forces the digit $6$, and it
would require $0 = \lfloor(3 \cdot 6 + c)/10\rfloor =
\lfloor(18 + c)/10\rfloor$ -- so only
$(7,2)$ survives in the interior: every digit is $7$,
$p = 77\ldots7$ and $q = 23\ldots32$. Now P1 follows
because even $m$ rules out the same-length case, which
needs $m$ odd; P3 from the alphabet union
$\{0,1,2\} \cup \{4,5,6\} \cup \{7\}$; P4 because a digit
$7$ excludes the same-length alphabet; and P2 because in
the non-repdigit same-length case $m$ is odd, so position
$m-2$ is odd: $a_{m-1} = a_0 \in \{1,2\}$ and
$a_{m-2} \in \{4,5,6\}$. For the image alphabet:
same-length case $b_j = 3a_j + 1 \in \{1,4,7\}$ at even
$j$ and $b_j = 3a_j \in \{12,15,18\}$ at odd $j$, digits
in $\{1,2,4,5,7,8\}$; overflow case $q = 23\ldots32$,
digits in $\{2,3\}$.
\end{proof}

\begin{remark}[P1 and the ``Includes A002281'' clause]
\label{rem:k3-iff}
P1 states the necessity direction; in the kernel the
even-length classification is the literal iff
\emph{even-length member $\Leftrightarrow$ repdigit
$77\ldots7$}, so Israel's ``Includes A002281'' clause --
every repdigit $77\ldots7$ is a member, image
$23\ldots32$ -- is itself certified
(Table~\ref{tab:lean}). P1--P4 are universal properties
of the members, not an enumeration of the sequence. The
$37$ published terms agree: their even-length members
are exactly $77$, $7777$, $777777$.
\end{remark}

% sec-A083831.tex -- k = 4 section of the unified kp+1
% paper. Standalone fragment, \input by
% A083833-UNIFIE-paper.tex.
% Environments and notation inherited from sec-carry.tex; the
% guards below are no-ops once defined there.
\providecommand{\IsPal}{\mathrm{IsPal}}
\makeatletter
\@ifundefined{theorem}{\newtheorem{theorem}{Theorem}}{}
\@ifundefined{corollary}{\newtheorem{corollary}[theorem]{Corollary}}{}
\makeatother

\section{Multiplier $k = 4$: A083831 and its image A083832}
\label{sec:k4}

Entry A083831~\cite{oeis} (A.~Murthy and M.~Srikanth,
May 2003; corrected and extended by R.~Zumkeller and
R.~Chandler) lists the palindromes $p$ such that
$q = 4p+1$ is again a palindrome:
$1, 2, 8, 88, 131, 141, 232, 242, 888, 8888, \ldots$
Commenting on the entry (April 4, 2018), Robert Israel
recorded the conjecture quoted in
Section~\ref{sec:intro}. We settle it for every length,
in both directions, together with the digit-alphabet
bound on $p$ and, as a corollary of the same carry
analysis, the alphabet of the companion sequence.

\begin{theorem}[Israel's conjecture on A083831, full form]
\label{thm:k4-even}
Let $m \ge 1$ be even and $\IsPal(p,m)$. Then $4p+1$ is
a palindrome if and only if $a_j = 8$ for every
$j < m$: the even-length members of A083831 are exactly
the repdigits $\underbrace{88\ldots8}_{m}$.
\end{theorem}

\begin{theorem}[digit alphabet of the source]
\label{thm:k4-dig}
Under the same hypotheses (any $m \ge 1$), every digit
of $p$ satisfies $a_j \le 4$ or $a_j = 8$.
\end{theorem}

\begin{corollary}[alphabet of the image, OEIS A083832]
\label{cor:k4-alphabet}
Every digit of $q = 4p+1$ lies in $\{1,2,3,5,6,9\}$.
\end{corollary}

\begin{proof}
Instantiating the shared carry chain at $k = 4$ gives
$c_j = \bigl\lfloor\bigl(4(p \bmod 10^j) +
1\bigr)/10^j\bigr\rfloor$, with $c_0 = 1$ and
$c_j \le 3$.
\begin{itemize}
\item \emph{Same length} ($m' = m$):
Lemma~\ref{lem:same} forces $m$ odd,
$a_0 \in \{1,2\}$, $a_j \in \{0,1,2\}$ at even positions
and $a_j \in \{3,4\}$ at odd positions. Output digits:
$4a_j + 1 \in \{1,5,9\}$ at even $j$ and
$4a_j \in \{12,16\}$ at odd $j$, i.e.\ digits
$\{2,6\}$.
\item \emph{Overflow} ($m' = m+1$). Since
$4a \bmod 10$ does not separate the carry, the direct
mirror law of the $k = 5$ case is unavailable; the
induction runs on the mirrored couple \emph{(A)}+\emph{(B)}
of Lemma~\ref{lem:overflow}, whose single fixed point is
$(a_{j+1}, c_{j+2}) = (8, 3)$: the chain is rigid in one
step. The head equation has the single solution
$a_0 = 8$ (checked digit by digit, the $m = 1$ boundary
$33 \bmod 10 = 3 = \lfloor 33/10 \rfloor$ included), so
$c_1 = 3$; propagating the fixed point, every digit is
$8$ and the carries freeze at $c_j = 3$:
$p = 88\ldots8$ and $q = 35\ldots53$ ($q = 33$ for
$m = 1$), digits in $\{3,5\}$.
\end{itemize}
For $m$ even the same-length case is impossible, so the
overflow case applies and the forward direction of
Theorem~\ref{thm:k4-even} follows; the converse
direction is the certified repdigit membership
(Table~\ref{tab:lean}). The union of the two digit
classes gives Theorem~\ref{thm:k4-dig}, and reading
$b_j$ in both cases gives
$\{1,2,5,6,9\} \cup \{3,5\} = \{1,2,3,5,6,9\}$,
Corollary~\ref{cor:k4-alphabet}.
\end{proof}

The published terms agree: the even-length members are
exactly $88$, $8888$, $888888$; the image family
$33, 353, 3553, 35553, \ldots$ is the overflow branch
and $525, 565, 929, 969, 52125, \ldots$ the same-length
one.

% sec-A083833.tex -- k = 5 section of the unified kp+1
% paper. Standalone fragment, \input by
% A083833-UNIFIE-paper.tex.
% Environments and notation inherited from sec-carry.tex; the
% guards below are no-ops once defined there.
\providecommand{\IsPal}{\mathrm{IsPal}}
\makeatletter
\@ifundefined{theorem}{\newtheorem{theorem}{Theorem}}{}
\@ifundefined{corollary}{\newtheorem{corollary}[theorem]{Corollary}}{}
\@ifundefined{remark}{\newtheorem{remark}[theorem]{Remark}}{}
\makeatother

\section{Multiplier $k = 5$: A083833 and its image A083834}
\label{sec:k5}

OEIS A083833~\cite{oeis} (Murthy and Srikanth, 2003)
lists the palindromes $p$ such that $q = 5p+1$ is again
a palindrome; A083834 is the image sequence of the
values $q$. On A083833, Michael S.~Branicky recorded
(June 13, 2021, tagged ``(Conjectures)''): ``All terms
contain only the digits $\{0,1,2,3\}$'' and ``For $d$
even, the only term with $d$ digits is $2^d$;
equivalently, all terms starting with $1$ have odd
length'' (in the entry's notation the caret $2^d$
denotes the $d$-fold concatenated repdigit $22\ldots2$).
On the image entry A083834 he recorded the analogous
pair, the alphabet $\{0,1,5,6\}$ and a repunit family;
that verbatim contains a parity slip discussed in
Remark~\ref{rem:parity-slip}.

\begin{theorem}[Branicky C1: alphabet of $p$]
\label{thm:k5-c1}
For all naturals $p, m, m'$ with $\IsPal(p,m)$,
$\IsPal(5p+1,m')$ and $m \ge 1$, every digit of $p$
satisfies $a_j \le 3$ for all $j < m$.
\end{theorem}

\begin{theorem}[Branicky C2: even length, full form]
\label{thm:k5-c2}
Let $m \ge 1$ be even and $\IsPal(p,m)$. Then $5p+1$ is
a palindrome if and only if $a_j = 2$ for all $j < m$:
the even-length members are exactly the repdigits
$p = 22\ldots2$.
\end{theorem}

A term starting with $1$ is not the repdigit, so by
Theorem~\ref{thm:k5-c2} it cannot have even length:
this is Branicky's ``equivalently'' clause, a
paper-side corollary of C2 rather than a separate
kernel declaration.

\begin{proof}[Proof of
Theorems~\ref{thm:k5-c1}--\ref{thm:k5-c2} and
Corollary~\ref{cor:k5-image}]
The shared carry-chain device (Section~\ref{sec:carry})
instantiates as follows. The carry into position $j$ is
$c_j = \bigl\lfloor\bigl(5\,(p \bmod 10^j)+1\bigr)
/10^j\bigr\rfloor$, bounded by $4$, with $c_0 = 1$; the
output length satisfies $m' \in \{m, m+1\}$. If
$m' = m$, Lemma~\ref{lem:same} forces $m$ odd,
$a_0 = 1$, $c_1 = 0$, and the alternating classes
$a_j \in \{0,1\}$ at even $j$, $a_j \in \{2,3\}$ at odd
$j$. If $m' = m+1$, the swap-freeness collapses the
chain to the constant $1$ (Scheme II of
Lemma~\ref{lem:overflow}); the boundary equations give
$a_0 = 2$, and the parity propagation
$a_j \equiv a_{j-1} \pmod 2$ pins every digit to $2$:
$p = 22\ldots2$ and $q = 11\ldots1$ is a repunit.
Reading $b_j = (5a_j + c_j) \bmod 10$ off the two
branches gives the image alphabet $\{1,6\}$ at even $j$,
$\{0,5\}$ at odd $j$, and the repunit $1\ldots1$ at
overflow.
Theorem~\ref{thm:k5-c1} is the union of the two digit
classes; if $m$ is even the same-length alternative is
excluded, which is the forward direction of
Theorem~\ref{thm:k5-c2}; the converse direction is the
certified identity
$5\cdot 2R_m + 1 = R_{m+1}$ packaged as a literal iff
(Table~\ref{tab:lean}). The $29$ listed terms fit:
their even-length members are exactly $22$, $2222$,
$222222$, $22222222$.
\end{proof}

\begin{remark}[parity slip in the A083834 verbatim]
\label{rem:parity-slip}
On A083834 the comment reads ``For $d$ \emph{odd}, the
only term with $d$ digits is $1^d$; equivalently, all
terms starting with $6$ have odd length''. The ``$d$
odd'' is a slip: the listed odd-length terms $606$ and
$656$ are not repunits, so the literal reading is
refuted by the entry's own data, whereas the
``equivalently'' clause is unambiguous. The corrected
statement is that even-length terms are repunits, and
that is the statement the kernel proves. On A083833
the analogous comment had been recorded correctly with
``$d$ even''.
\end{remark}

\begin{corollary}[image alphabet: A083834]
\label{cor:k5-image}
Under the same hypotheses, every digit of $q = 5p+1$
lies in $\{0,1,5,6\}$; if $m'$ is even then every digit
of $q$ equals $1$, i.e.\ $q$ is the repunit
$11\ldots1$; and a term of A083834 starting with $6$
has odd length.
\end{corollary}

% sec-A083835.tex -- k = 6 section of the unified kp+1
% paper. Standalone fragment, \input by
% A083833-UNIFIE-paper.tex.
% Environments and notation inherited from sec-carry.tex; the
% guards below are no-ops once defined there.
\providecommand{\IsPal}{\mathrm{IsPal}}
\makeatletter
\@ifundefined{theorem}{\newtheorem{theorem}{Theorem}}{}
\@ifundefined{corollary}{\newtheorem{corollary}[theorem]{Corollary}}{}
\@ifundefined{remark}{\newtheorem{remark}[theorem]{Remark}}{}
\makeatother

\section{Multiplier $k = 6$: A083835 and its image A083836}
\label{sec:k6}

OEIS A083835~\cite{oeis} (Murthy and Srikanth, May 2003;
extended by R.~Chandler) lists the palindromes $p$ such
that $q = 6p+1$ is again a palindrome:
$1$, $9$, $99$, $121$, $131$, $999$, $9999$, $12021$,
$12121$, $13031$, $13131$, \ldots; A083836 is the image
sequence $7$, $55$, $595$, $727$, $787$, $5995$,
$59995$, $72127$, \ldots
No conjecture is recorded on either entry. At this
multiplier the shared device yields not only the
necessary structure but the complete classification,
proved in both directions in the kernel.

\begin{theorem}[complete classification of A083835]
\label{thm:k6-iff}
Let $\IsPal(p,m)$ with $m \ge 1$. Then $6p+1$ is a
palindrome of some length $m'$ if and only if exactly
one of the following holds:
\begin{itemize}
\item[(i)] $m' = m + 1$ and every digit of $p$ equals
$9$, i.e.\ $p = 99\ldots9$ is the repdigit $9\,R_m$;
\item[(ii)] $m' = m$, $m$ is odd, and the digits of $p$
obey the alternating classes $a_j \in \{0,1\}$ at even
positions and $a_j \in \{2,3\}$ at odd positions.
\end{itemize}
\end{theorem}

\begin{corollary}[image companions: A083836]
\label{cor:k6-img}
Under the same hypotheses, every digit of $q = 6p+1$
avoids $\{0,3,4,6\}$, i.e.\ lies in $\{1,2,5,7,8,9\}$.
If $m'$ is even then $q$ has the exact form
$5\,9^{\,m'-2}\,5$: endpoints $5$, all interior digits
$9$ (the images of the repdigit branch).
\end{corollary}

\begin{proof}
Necessity is the instantiation of
Section~\ref{sec:carry} at $k = 6$: carries
$c_j = \bigl\lfloor\bigl(6\,(p \bmod 10^j)+1\bigr)
/10^j\bigr\rfloor \le 5$ with $c_0 = 1$. At constant
length, Lemma~\ref{lem:same} gives $m$ odd, $a_0 = 1$
and the classes $\{0,1\}$/$\{2,3\}$; at overflow, the
head equation $(6a_0+1) \bmod 10 = c_m \in \{1,\dots,5\}$
has the single solution $a_0 = 9$, $c_m = 5$, and the
mirrored local system \emph{(A)}+\emph{(B)} of
Lemma~\ref{lem:overflow} has the unique fixed point
$(a',c') = (9,5)$ -- verified by a finite check over the
ten digits, with no residual spurious state (unlike
$k = 3$): every digit is $9$ and
$q = 5\,9^{\,m-1}\,5$. For the converse directions the
carry chain is rebuilt from the prescribed digits: an
all-$9$ palindrome of length $m$ gives
$6p+1 = 5\,9^{\,m-1}\,5$, a palindrome of length $m+1$;
an odd-length palindrome with the alternating classes
closes the chain $c_j = 1$ at even $j$, $c_j = 0$ at odd
$j$, gives $c_m = 0$ and a palindromic $6p+1$ of length
$m$. Both directions are kernel-verified
(Table~\ref{tab:lean}). The image alphabet reads off
the branches: same-length digits in $\{1,2,7,8\}$,
overflow digits in $\{5,9\}$.
\end{proof}

\begin{remark}[the even-length members of A083835]
\label{rem:k6-even}
Since the same-length branch requires $m$ odd, every
even-length member of A083835 is the repdigit
$99\ldots9$, with image $5\,9^{\,m'-2}\,5$; the
published even-length terms are exactly $99$, $9999$,
$999999$, $99999999$. The same parity argument gives
$q$ even-length only in the overflow branch, which is
the second clause of Corollary~\ref{cor:k6-img}.
\end{remark}

\begin{remark}[the shared alternating classes at
$k = 5$ and $k = 6$]
\label{rem:converse}
Table~\ref{tab:classes} shows the same-length classes
coincide at $k = 5$ and $k = 6$: $\{0,1\}$ at even
positions, $\{2,3\}$ at odd positions, $a_0 = 1$. Hence
every non-repdigit member of A083833 carries a digit
structure which Theorem~\ref{thm:k6-iff} certifies to
be a member of A083835: the inclusion
\[
\{\text{non-repdigit terms of A083833}\}
\;\subseteq\;
\{\text{non-repdigit terms of A083835}\}
\]
is kernel-verified, and the $29$ published terms of the
two entries display the same non-repdigit members
($121, 131, 12021, 12121, 13031, \ldots$). The reverse
inclusion follows on paper: for the same classes the
$k = 5$ chain closes with the same alternating carries
$1,0,1,0,\ldots$ (the generic instance is
Lemma~\ref{lem:altsuff}) -- the exact analogue of the certified
$k = 6$ converse, and likewise kernel-certified at
$k = 5$ by \texttt{alt\_gives\_pal}
(Remark~\ref{rem:sufftask}).
\end{remark}

% sec-zeros.tex -- the certified zeros of A083837 and
% A083838. Standalone fragment, \input by
% A083833-UNIFIE-paper.tex.
% Environments and notation inherited from sec-carry.tex; the
% guards below are no-ops once defined there.
\providecommand{\revW}{\mathrm{rev}}
\makeatletter
\@ifundefined{theorem}{\newtheorem{theorem}{Theorem}}{}
\@ifundefined{lemma}{\newtheorem{lemma}[theorem]{Lemma}}{}
\@ifundefined{corollary}{\newtheorem{corollary}[theorem]{Corollary}}{}
\@ifundefined{remark}{\newtheorem{remark}[theorem]{Remark}}{}
\makeatother

\section{The certified zeros of A083837 and A083838}
\label{sec:zeros}

The last two entries of the family scan the multiplier.
A083837~\cite{oeis} (Murthy and Srikanth, May 2003;
corrected and extended by R.~Chandler) is quoted in
Section~\ref{sec:intro}: $a(k)$ is the smallest
palindrome $p$ such that $kp+1$ is palindromic, or $0$
if none exists. Its $104$ displayed terms vanish at the
$21$ indices
\[
 k \in Z = \{12,17,23,31,37,41,42,48,51,53,61,64,71,
 74,79,81,86,91,96,101,103\},
\]
where Chandler's comment reports only the absence of a
palindrome below $10^{15}$. A displayed zero is thus
finite evidence, not nonexistence; a later witness could
refute it. This section upgrades each of the $21$ zeros
to a universal theorem, by an obstruction on the two
ends of the candidate palindrome.

\paragraph{The prefix-suffix obstruction.}
Let $p$ be an $m$-digit palindrome, $1 \le w \le m$, and
let $P$ be its $w$-prefix (the integer formed by its
first $w$ digits). Write $\revW_w(n)$ for the integer
read from the last $w$ digits of $n$ in reverse order.
Assume $q = kp+1$ is a palindrome of $m'$ digits. Then:
\begin{itemize}
\item the prefix bounds the value:
$kP\,10^{m-w} \le q < k(P+1)\,10^{m-w}$;
\item the last $w$ digits of $p$ form $\revW_w(P)$
(palindromicity of $p$), i.e.\
$p \equiv \revW_w(P) \pmod{10^w}$, so the last $w$
digits of $q$ are those of $k\,\revW_w(P) + 1$; as $q$ is
itself palindromic of length $m' \ge m \ge w$, its
$w$-prefix is the reversal of its last $w$ digits:
$L = \revW_w\bigl(k\,\revW_w(P) + 1\bigr)$. If
$L < 10^{w-1}$ the last digit of $q$, hence its leading
digit, would be $0$ -- impossible, so $L \ge 10^{w-1}$
and $L\,10^{m'-w} \le q < (L+1)\,10^{m'-w}$;
\item writing $s = m' - m \ge 0$, one has $s \le S$ for
any $S$ with $k < 10^S$, since
$q < k\,10^m + 1 \le 10^{m+S}$; crossing the two windows
($L\,10^{m'-w} \le q < k(P+1)\,10^{m-w}$ and
$kP\,10^{m-w} \le q < (L+1)\,10^{m'-w}$) and dividing by
$10^{m-w}$ then forces
\begin{equation}
 L\,10^{s} \;\le\; k\,(P+1)
 \qquad\text{and}\qquad
 k\,P \;<\; (L+1)\,10^{s} .
 \label{eq:overlap}
\end{equation}
\end{itemize}

\begin{lemma}[border obstruction]
\label{lem:border}
If, for a candidate $w$-prefix $P$, either
$L < 10^{w-1}$ (the required $w$-prefix of $q$ would
start with $0$, i.e.\ $q$ would end in $0$) or no
$s \in \{0,\dots,S\}$ satisfies \eqref{eq:overlap},
then no $m$-digit palindrome $p$ with $w$-prefix $P$
can have $kp+1$ palindromic. In particular, if every
$w$-prefix $P \in \{10^{w-1},\dots,10^{w}-1\}$ is
obstructed, no palindrome $p$ of length $\ge w$ has
$kp+1$ palindromic.
\end{lemma}

The lemma prunes the prefix tree level by level: the
survivors of depth $w$ seed the candidate $(w+1)$-prefixes
$10P + d$ for $d \in \{0,\dots,9\}$. The residual cases
$m < w$ are finite: an $m$-digit palindrome is its own
$m$-prefix, so it must appear among the level-$m$
survivors, be palindromic as a prefix, and pass the
direct check on $kp+1$. Table~\ref{tab:zeros} gives the
certificate tree for each of the $21$ indices: the depth
$w$ at which the survivor list becomes empty, and the
survivors $E_t$ at each shallower level.

\begin{table}[ht]
\centering
\small
{\footnotesize\setlength{\tabcolsep}{4pt}
\begin{tabular}{>{\raggedright\arraybackslash}p{4.4cm}|c|%
>{\raggedright\arraybackslash}p{10.2cm}}
\hline
$k$ & depth & survivors \\
\hline
$12$, $31$, $37$, $41$, $42$, $51$, $53$, $61$, $64$,
$71$, $81$, $86$, $91$, $96$ & $1$ & none: every leading
digit $a_0 \in \{1,\dots,9\}$ is obstructed \\
$17$ & $2$ & $E_1 = \{2\}$; $m = 1$: $17\cdot2+1 = 35$
not palindromic; $E_2 = \varnothing$ \\
$23$ & $2$ & $E_1 = \{1\}$; $m = 1$: $23\cdot1+1 = 24$
not palindromic; $E_2 = \varnothing$ \\
$74$ & $2$ & $E_1 = \{6,9\}$; $m = 1$: $445$, $667$
not palindromic; $E_2 = \varnothing$ \\
$101$ & $2$ & $E_1 = \{1,\dots,8\}$ ($S = 3$);
$m = 1$ excluded case by case; $E_2 = \varnothing$ \\
$103$ & $2$ & $E_1 = \{5\}$ ($S = 3$); $m = 1$:
$516$ not palindromic; $E_2 = \varnothing$ \\
$48$ & $3$ & $E_1 = \{1\}$; $m = 1$: $48\cdot1+1 = 49$
not palindromic; $E_2 = \{19\}$; $m = 2$: candidate
prefix $19$ not a palindrome ($9 \ne 1$);
$E_3 = \varnothing$ over $\{190,\dots,199\}$ \\
$79$ & $3$ & $E_1 = \{6\}$; $m = 1$: $79\cdot6+1 = 475$
not palindromic; $E_2 = \{65\}$; $m = 2$: candidate
prefix $65$ not a palindrome; $E_3 = \varnothing$
over $\{650,\dots,659\}$ \\
\hline
\end{tabular}}
\caption{Certificate trees of the $21$ certified zeros
of A083837. $E_t$ is the list of $t$-prefixes escaping
the obstruction at level $t$; $m < w$ cases are
finished by direct products.}
\label{tab:zeros}
\end{table}

\begin{theorem}[the $21$ zeros of A083837]
\label{thm:zeros}
For every $k \in Z$ there is no palindrome $p$ of any
length such that $kp+1$ is a palindrome.
\end{theorem}

\begin{proof}
For each $k$, the obstruction of
Lemma~\ref{lem:border} is evaluated on the candidate
prefixes level by level: at depth $1$ the list
$E_1 = \varnothing$ for the fourteen indices of the
first row of Table~\ref{tab:zeros}; at depth $2$ for
$k \in \{17,23,74,101,103\}$ the level-$1$ survivors
shown are excluded at length $m = 1$ by the displayed
products and the next level is empty; at depth $3$ for
$k \in \{48,79\}$ the chains shown close the tree.
Every case is kernel-checked in \texttt{A083837.lean}
(Table~\ref{tab:lean}).
\end{proof}

\begin{corollary}[the $21$ zeros of A083838]
\label{cor:zeros838}
For every $k \in Z$ there is no pair of palindromes
$(p,q)$ with $q = kp+1$: the image entry A083838
vanishes at the same indices.
\end{corollary}

The companion entry A083838 displays $66$ terms; its
$12$ shown zeros, at indices $12$, $17$, $23$, $31$,
$37$, $41$, $42$, $48$, $51$, $53$, $61$ and $64$, are
exactly the displayed part of the certified set, while
the remaining nine indices $71$, $74$, $79$, $81$,
$86$, $91$, $96$, $101$ and $103$ lie beyond its
displayed window and are settled here on the same
footing.

\begin{proof}
Each corollary is the direct negation of the existence
of a member pair, a one-line consequence of
Theorem~\ref{thm:zeros}; the $21$ statements are
kernel-checked in \texttt{A083838.lean}.
\end{proof}

\begin{remark}[scope of the certification]
\label{rem:scopezeros}
Theorem~\ref{thm:zeros} is a universal nonexistence
statement: it rules out every palindrome of every
length, where the displayed zero reported only
$p < 10^{15}$. It says nothing about the nonzero
displayed values of A083837 or A083838, which are not
certified here; and it does not address multipliers
$k > 104$, beyond the displayed range of the entry.
Consistently, an independent enumeration found no
witness palindrome below $10^{11}$ for any of the $21$
indices -- agreement with the finite-search comment,
though never used as evidence.
\end{remark}


\section{Formal verification}
\label{sec:lean}

It is important to distinguish what is kernel-verified in
Lean~4 from what belongs to the mathematical presentation
of this paper:
\begin{itemize}
\item \emph{Kernel-verified in Lean~4~\cite{lean4}
(v4.35.0) on top of mathlib~\cite{mathlib}
(v4.35.0)~\cite{artifacts}:} every stated theorem, lemma
and corollary of Sections~\ref{sec:carry}--\ref{sec:zeros}
-- the per-multiplier carry-chain machinery, the five
classification modules \texttt{A069881.lean},
\texttt{A083829.lean}, \texttt{A083831.lean},
\texttt{A083833.lean} and \texttt{A083835.lean} (the
last containing the complete iff and the three A083836
image companions), the thin image re-export modules
\texttt{A083830.lean}, \texttt{A083832.lean},
\texttt{A083834.lean} and \texttt{A083836.lean}, the
converse module \texttt{Repdigits\_KP1.lean} holding the
four literal repdigit iffs, and the obstruction module
\texttt{A083837.lean} with the image corollary module
\texttt{A083838.lean}. A deposited specification module
per entry closes each capstone against a frozen
specification, and sanity modules anchor the
statements on published terms and certify that each
proved class is inhabited.
\item \emph{Proved in this paper only:} the repdigit
families of Remark~\ref{rem:k1} for $k \in \{7,8,9\}$
and the elementary $k = 1$ remark. Neither is stated as
a kernel theorem.
\end{itemize}

Table~\ref{tab:lean} records the correspondence between
the statements above and the Lean declarations. The
formalization is duplicated per multiplier
(Remark~\ref{rem:honest}): each module carries its own
copy of the carry chain, so the correspondence is
declared per $k$.

\begin{table}[ht]
\centering
\small
{\footnotesize\setlength{\tabcolsep}{4pt}
\begin{tabular}{>{\raggedright\arraybackslash}p{4.4cm}%
>{\raggedright\arraybackslash}p{6.0cm}%
>{\raggedright\arraybackslash}p{3.8cm}}
\hline
Paper item & Lean declaration & Module \\
\hline
Lemma~\ref{lem:same} ($k=2$) &
\texttt{crry\_sym\_same}, \texttt{carry\_period\_same},
\texttt{same\_case\_odd\_a0}, \texttt{same\_case\_digits} &
\texttt{A069881.lean} \\
Theorem~\ref{thm:k2-even} & \texttt{A069881\_even} (forward),
\texttt{repdigit5\_iff\_member\_even\_}\allowbreak
\texttt{A069881} (iff) &
\texttt{A069881.lean}, \texttt{Repdigits\_KP1.lean} \\
Theorem~\ref{thm:k2-dicho}, Cor.~\ref{cor:k2-img} &
\texttt{A069881\_overflow}, \texttt{succ\_case\_all5},
\texttt{A069881\_image\_odd} & \texttt{A069881.lean} \\
overflow iff, $k=2$ &
\texttt{repdigit5\_iff\_overflow\_A069881} &
\texttt{Repdigits\_KP1.lean} \\
Theorem~\ref{thm:k3-israel} (P1--P4) &
\texttt{A083829\_even}, \texttt{A083829\_prefix},
\texttt{A083829\_digits}, \texttt{A083829\_seven} &
\texttt{A083829.lean} \\
P1 as iff & \texttt{repdigit7\_iff\_member\_even\_}%
\allowbreak \texttt{A083829} &
\texttt{Repdigits\_KP1.lean} \\
Corollary~\ref{cor:k3-alphabet} &
\texttt{A083830\_digits}, re-export
\texttt{A083830\_alphabet} & \texttt{A083829.lean},
\texttt{A083830.lean} \\
Theorem~\ref{thm:k4-even} & \texttt{A083831\_even}
(forward), \texttt{repdigit8\_iff\_member\_even\_}%
\allowbreak \texttt{A083831} (iff) & \texttt{A083831.lean},
\texttt{Repdigits\_KP1.lean} \\
Theorem~\ref{thm:k4-dig} & \texttt{A083831\_digits} &
\texttt{A083831.lean} \\
Corollary~\ref{cor:k4-alphabet} &
\texttt{A083832\_digits}, re-export
\texttt{A083832\_alphabet} & \texttt{A083831.lean},
\texttt{A083832.lean} \\
Theorem~\ref{thm:k5-c1}, Theorem~\ref{thm:k5-c2} &
\texttt{A083833\_C1}, \texttt{A083833\_C2} &
\texttt{A083833.lean} \\
C2 as iff & \texttt{repdigit2\_iff\_member\_even\_}%
\allowbreak \texttt{A083833} &
\texttt{Repdigits\_KP1.lean} \\
Corollary~\ref{cor:k5-image} &
\texttt{A083834\_C1}, \texttt{A083834\_C2}, re-exports
\texttt{A083834\_alphabet}, \texttt{A083834\_repunit},
\texttt{A083834\_starts6\_odd} & \texttt{A083833.lean},
\texttt{A083834.lean} \\
Theorem~\ref{thm:k6-iff} & \texttt{A083835} (complete iff;
directions \texttt{all9\_gives\_pal},
\texttt{alt\_gives\_pal}) & \texttt{A083835.lean} \\
Corollary~\ref{cor:k6-img} &
\texttt{A083836\_digits}, \texttt{A083836\_even},
\texttt{A083836\_form}, re-exports
\texttt{A083836\_alphabet}, \texttt{A083836\_evenlen},
\texttt{A083836\_shape} & \texttt{A083835.lean},
\texttt{A083836.lean} \\
Lemma~\ref{lem:altsuff} & \texttt{alt\_gives\_pal}
(one instance per multiplier, $k \in \{2,3,4,5,6\}$) &
\texttt{A069881.lean}, \texttt{A083829.lean},
\texttt{A083831.lean}, \texttt{A083833.lean},
\texttt{A083835.lean} \\
Lemma~\ref{lem:border}, Theorem~\ref{thm:zeros} &
\texttt{revW}, \texttt{okP}, \texttt{escapesIn},
\texttt{prefix\_escapes}, \texttt{zero\_of\_escapes},
\texttt{no\_wit\_}$k$, \texttt{A083837\_}$k$\texttt{\_zero}
($21$ indices $k$) & \texttt{A083837.lean} \\
Corollary~\ref{cor:zeros838} &
\texttt{A083838\_}$k$\texttt{\_zero} ($21$ corollaries) &
\texttt{A083838.lean} \\
\hline
\end{tabular}}
\caption{Correspondence between mathematical statements
and Lean~4 declarations. The $21$ zero indices are
$\{12,17,23,31,37,41,42,48,51,53,61,64,71,74,79,81,86,
91,96,101,103\}$.}
\label{tab:lean}
\end{table}

The axiom audits report, for every declaration listed in
Table~\ref{tab:lean}, only the three standard axioms
\texttt{propext}, \texttt{Classical.choice} and
\texttt{Quot.sound}; a source scan confirms zero
\texttt{sorry} and zero \texttt{native\_decide}. Finite
digit checks (small cases, the per-$k$ head equation,
the escape lists of Section~\ref{sec:zeros}) are decided
inside the kernel by \texttt{decide} and
\texttt{interval\_cases} on closed values; the carry-chain
and obstruction arguments are ordinary proofs. The sealed
version of the A083833 module was additionally replayed
declaration-by-declaration through three external kernel
checkers~\cite{kernels} -- \texttt{nanoda},
\texttt{con-leche} and its Rust port \texttt{con-ron} --
each at exit status $0$; the other modules rely on the
kernel build and the axiom audit, the same trust base as
their previously deposited artifacts~\cite{artifacts}.

\paragraph{Scope.} The certified theorems are universal
properties of the members -- digit classes, parity of the
length, repdigit and repunit families, the complete
classification at $k = 6$, and universal nonexistence at
the twenty-one zero indices. They are consistent with
every published term of the eleven entries; completeness
of the displayed lists is a separate empirical matter,
not claimed here. The nonzero displayed values of
A083837 and A083838 are likewise outside the certified
perimeter: what is proved is the zero clause, for every
candidate palindrome of any length. Convention: the
membership predicate is taken over $m \ge 1$ digit
expansions without leading zero, so $p = 0$ is outside
the specification domain throughout.

\begin{remark}[the converse is kernel-verified at every
$k$]
\label{rem:sufftask}
At $k = 6$ the kernel proves the complete iff
(Theorem~\ref{thm:k6-iff}): the alternating-class
condition is both necessary and sufficient.
Lemma~\ref{lem:altsuff} is its paper-level umbrella
across multipliers, and the kernel now carries the same
theorem inside each classification module as
\texttt{alt\_gives\_pal} (Table~\ref{tab:lean}).
\end{remark}

\section{Open items}
\label{sec:open}

\begin{question}[the family beyond $k = 6$]
\label{q:beyond}
The head key $kd+1 = 11v$ does not by itself select a
repdigit overflow family: for $p = d\,R_m$ one has
$kp+1 = v\,10^m + (2v-1)\,(10^m-10)/9 + v$, the digit
string $v\,(2v-1)^{m-1}\,v$, which is a palindrome
without carries exactly when $2v-1 \le 9$, i.e.\
$v \le 5$. Beyond the certified perimeter the digit
solutions $(d,v) = (1,1)$ at $k = 10$, $(2,3)$ at
$k = 16$ and $(3,5)$ at $k = 18$ satisfy the bound and
give the infinite families $1^m \mapsto 1^{m+1}$,
$2^m \mapsto 3\,5^{\,m-1}\,3$ and
$3^m \mapsto 5\,9^{\,m-1}\,5$, matching the displayed
$a(10) = 1$, $a(16) = 2$, $a(18) = 3$; the solutions
$(5,6)$ at $k = 13$ and $(7,9)$ at $k = 14$ fail it --
$13\cdot55+1 = 716$ and $14\cdot77+1 = 1079$ are not
palindromes -- and the repdigit branch yields only the
displayed one-digit witnesses $a(13) = 5$, $a(14) = 7$
(at $k = 17$ the congruence forces $v = 14$, off the
digit range, consistent with the certified zero);
witnesses also occur without any head-key solution, as
$a(11) = 55$ through $11\cdot55+1 = 606$. For
$k \ge 10$ the bound $c_j \le k-1$ leaves the
single-digit regime and the carry device needs
extension. Whether the sets $S_k$ for
$k \in \{7,8,9\}$ -- and beyond -- admit the same
two-branch classification, and whether the certified
zero set exhausts the zero set of the entry, is open.
\end{question}

\section{Code availability}
\label{sec:code}

The Lean~4 sources of the previously deposited A083833
development are archived in the open
record~\cite{artifacts}; the consolidated development
covering the eleven sequences treated here -- the five
classification modules, the image re-exports, the
converse module \texttt{Repdigits\_KP1.lean} and the
obstruction modules \texttt{A083837.lean} and
\texttt{A083838.lean}, together with the specification
and sanity modules and the pinned build files -- is
deposited alongside this article in its verification
record.

\section*{Use of AI tools}
Generative AI tools assisted in OEIS transcription,
auditing, proof sketching, and text review, within a
verification methodology designed and supervised by the
author. The theorems of
Sections~\ref{sec:carry}--\ref{sec:zeros} are verified
from standard axioms only in the Lean~4 kernel, with the
Lean source archived in the verification record. The
author selected the problems and proof strategies,
arbitrated all discrepancies, verified OEIS fidelity,
and takes full responsibility for this work.

\section*{Funding}
No funding was received for this work.

\section*{Competing interests}
The author declares none.

\begin{thebibliography}{99}

\bibitem{oeis}
OEIS Foundation Inc.~(2026), \emph{The On-Line
Encyclopedia of Integer Sequences},
\url{https://oeis.org}; entries A069881 (A.~Murthy,
April 30 2002), A083829--A083838 (A.~Murthy and
M.~Srikanth, May 9 2003), all corrected and extended by
R.~Chandler, May 2003.

\bibitem{tangente}
\emph{Tangente}, problem~8154 (the only even-digit
palindromes $\alpha+$ are $55, 5555, 555555, \ldots$;
as cited on OEIS A069881).

\bibitem{a082274}
OEIS Foundation Inc.~(2026), \emph{Entry A082274 in The
On-Line Encyclopedia of Integer Sequences},
\url{https://oeis.org/A082274} (A.~Murthy, April 13
2003; characterization proved on the entry by
F.~Fr\"ohlich, July 2014).

\bibitem{a082275}
OEIS Foundation Inc.~(2026), \emph{Entry A082275 in The
On-Line Encyclopedia of Integer Sequences},
\url{https://oeis.org/A082275} (A.~Murthy, April 13
2003; characterization proved on the entry by
D.~A.~Corneth, July 2025).

\bibitem{harmincsotak}
M.~Harminc and R.~Sot\'ak, \emph{Palindromic numbers in
arithmetic progressions}, Fibonacci Quart.\ \textbf{36}
(1998), no.~3, 259--262,
\url{https://doi.org/10.1080/00150517.1998.12428937}.

\bibitem{banks}
W.~D. Banks, D.~Hart and M.~Sakata, \emph{Almost all
palindromes are composite}, Math.\ Res.\ Lett.\
\textbf{11} (2004), 853--868,
\url{https://doi.org/10.4310/MRL.2004.v11.n6.a10}.

\bibitem{cll}
J.~Cilleruelo, F.~Luca and L.~Baxter, \emph{Every
positive integer is a sum of three palindromes}, Math.\
Comp.\ \textbf{87} (2018), 3023--3055,
\url{https://doi.org/10.1090/mcom/3221}.

\bibitem{holte}
J.~M. Holte, \emph{Carries, combinatorics, and an
amazing matrix}, Amer.\ Math.\ Monthly \textbf{104}
(1997), no.~2, 138--149,
\url{https://doi.org/10.1080/00029890.1997.11990612}.

\bibitem{lean4}
L.~de~Moura and S.~Ullrich, \emph{The Lean 4 theorem
prover and programming language}, in A.~Platzer and
G.~Sutcliffe (eds.), CADE-28, Lecture Notes in
Comput.\ Sci.\ 12699 (2021) 625--635.

\bibitem{mathlib}
The mathlib Community, \emph{The Lean mathematical
library}, CPP '20: Proc.\ 9th ACM SIGPLAN Int.\ Conf.\
Certified Programs and Proofs (2020) 367--381.

\bibitem{artifacts}
F.~Gaudin, \emph{Consolidated Lean~4 formalization of
the $kp+1$ palindrome family (OEIS A069881, A083829,
A083830, A083831, A083832, A083833, A083834, A083835,
A083836, A083837, A083838)}, Lean~4 + mathlib
(v4.35.0); package deposited alongside this article
\url{https://doi.org/10.5281/zenodo.24420836}.
The previously deposited A083833
development is archived at
\url{https://doi.org/10.5281/zenodo.24078012} (open).

\bibitem{kernels}
Independent third-party Lean~4 kernel checkers used for
the export replay (separate projects, not part of the
\texttt{leanprover/lean4:v4.35.0} distribution, which
ships its own \texttt{leanchecker}): \texttt{nanoda\_bin}
(Rust kernel checker), C.~Bailey,
\url{https://github.com/ammkrn/nanoda_lib}, commit
\texttt{0da9f98}; \texttt{con-leche} (Lean~4 kernel
checker implemented in Lean, with a machine-checked
consistency proof), Lean FRO,
\url{https://github.com/leanprover/con-leche}, commit
\texttt{d7d2b7c}; \texttt{con-ron} (Rust port of the
\texttt{con-leche} kernel checker), Lean FRO,
\url{https://github.com/leanprover/con-ron}, commit
\texttt{42a0864}. Commits pinned as of 2026-11-13.

\end{thebibliography}

\end{document}
